% !TeX program = pdflatex
% English release manuscript; the local working draft is maintained separately.
% Official NeurIPS 2026 style:
% https://media.neurips.cc/Conferences/NeurIPS2026/Formatting_Instructions_For_NeurIPS_2026.zip
\documentclass{article}
\PassOptionsToPackage{numbers,sort&compress}{natbib}
\usepackage[preprint]{neurips_2026}
\usepackage[T1]{fontenc}
\usepackage[utf8]{inputenc}
\usepackage{amsmath,amssymb}
\usepackage{booktabs,tabularx,array}
\usepackage{microtype}
\usepackage[hidelinks]{hyperref}
\newcolumntype{Y}{>{\raggedleft\arraybackslash}X}
\title{TAPE: Task-Aware Evolutionary Expert Pruning under a Fixed Global Budget}
\author{DSAI6707 Project Report}
\hypersetup{pdftitle={TAPE: Task-Aware Evolutionary Expert Pruning under a Fixed Global Budget}}
\begin{document}
\maketitle

\begin{abstract}
Expert pruning reduces the parameter footprint of sparse mixture-of-experts models. Its effectiveness depends on both layer-wise capacity allocation and the identities of retained experts. Existing importance-based methods commonly rank experts within each layer while fixing layer capacities in advance or optimizing them separately. We propose TAPE, a task-aware evolutionary pruning method that jointly searches layer capacities and expert identities under an exact global expert budget. Budget-preserving operators generate feasible candidates, and task-specific objectives combine correct-answer cross-entropy with agreement to the original model's answer distribution. No model weights are updated. Experiments on ERNIE-4.5-21B-A3B-PT and Granite-3.1-1B-A400M-Base use ARC-Easy and OpenBookQA to evaluate task performance after compression. TAPE achieves higher average accuracy than the direct-scoring baselines at both budgets on both models. At 50\% expert retention on Granite, it exceeds the combined-scoring baseline by 0.38 percentage points. Ablations support joint capacity and identity optimization; on ERNIE at 75\% retention, continued evolution improves average test accuracy by 1.57 percentage points over selection from the initial population.
\end{abstract}

\noindent\textbf{Keywords:} mixture of experts; expert pruning; evolutionary computation; multi-objective optimization; fixed budget; task awareness

\section{Introduction}
\label{sec:1}
Mixture-of-experts (MoE) models route inputs to different experts, expanding model capacity while limiting activated computation. Sparse routing reduces computation per input, but all expert weights must still be stored, making GPU memory a deployment constraint~\citep{ref1,ref3}. Expert pruning addresses this constraint by removing experts and reducing the model's parameter footprint~\citep{ref2,ref1}.

Pruning involves two coupled decisions: how many experts to retain in each layer, and which experts receive those slots. Uniform layer-wise pruning fixes capacities before selecting experts by importance. Non-uniform allocation permits different capacities, but a fixed within-layer ranking still restricts the expert sets that can be selected~\citep{ref17}. Separately optimizing capacity and identity does not explore their full joint space.

Expert importance is usually estimated from routing frequency, activation strength, or reconstruction error on calibration data~\citep{ref1,ref3}. These signals describe contributions on the calibration distribution, but do not directly evaluate the target-task performance of a complete retained set. Task-specific pruning already uses target-task data to estimate importance~\citep{ref2,ref3}; however, independently ranking experts differs from comparing complete expert combinations. Task information can guide the selection of those combinations as well as their evaluation.

TAPE, Task-Aware Pruning via Evolution, combines these two directions. It jointly searches layer capacities and expert identities at an identical global expert count, evaluating each retained set through prediction losses on the target question-answering tasks. Capacity/identity ablations, a before-and-after evolution comparison, and expert-selection analysis examine the resulting task gains.

Our contributions are:
\begin{enumerate}
\item A joint representation of layer capacities and expert identities under an exact global expert budget, allowing equal-budget comparisons of complete expert sets.
\item Budget-preserving identity swaps, cross-layer capacity transfers, and set crossover, combined with task-level CE/KL objectives and NSGA-II optimization.
\item Empirical analysis of equal-capacity expert selection, continued evolution, and expert-retention decisions relative to importance rankings.
\end{enumerate}

\section{Related work}
\label{sec:2}
\subsection{Expert scoring and compression}
\label{sec:2.1}
Importance-based compression evaluates output contributions, usage frequency, or functional similarity. Not All Experts are Equal~\citep{ref2} studies task-agnostic and task-specific pruning. EAC-MoE~\citep{ref3} uses expert-selection frequency, REAP~\citep{ref1} combines gate weights with expert-output norms, and STUN~\citep{ref4} combines expert-level and unstructured pruning. Direct ranking is inexpensive, but does not explicitly compare the downstream performance of complete pruned expert sets.

Expert merging and reconstruction provide alternatives. Merge, Then Compress~\citep{ref5}, HC-SMoE~\citep{ref6}, and Diversifying the Expert Knowledge~\citep{ref7} use alignment, output clustering, and similarity-based grouping to integrate expert knowledge. DERN~\citep{ref8} recombines neuron fragments from removed experts. CAMEx~\citep{ref9} and NAMEx~\citep{ref10} study curvature-aware merging and expert cooperation. TAPE keeps the expert weights unchanged and evaluates the selection of existing experts.

Other compression variables include low-rank sharing in MoE-SVD~\citep{ref11}, shared bases and expert differences in D\textsuperscript{2}-MoE~\citep{ref12}, shared subspaces in Sub-MoE~\citep{ref13}, and mixed precision and dynamic pruning in Mixture Compressor~\citep{ref14}. Our experiments instead fix an expert-parameter budget and examine expert-set selection.

\subsection{Evolutionary search and joint expert selection}
\label{sec:2.2}
EEP~\citep{ref15} applies gradient-free evolution to expert selection and extends its representation to routing maps and expert-merging matrices. MoE-I\textsuperscript{2}~\citep{ref16} first determines layer-wise pruning amounts, then performs within-layer genetic search and small-block combination; its full pipeline also includes low-rank decomposition and recovery fine-tuning. EvoESAP~\citep{ref17} searches layer capacities under a fixed global budget while retaining fixed within-layer expert rankings.

Evolutionary Optimization of Model Merging Recipes~\citep{ref18}, CycleQD~\citep{ref19}, and MERGE\textsuperscript{3}~\citep{ref20} study combination representations, task diversity, and evolutionary optimization with limited evaluation budgets. Together, these works motivate attention to both the search representation and its evaluation objectives.

TAPE encodes capacities and identities across all MoE layers in one candidate. Unlike EEP's routing and merging matrices, it directly selects a binary subset of existing experts, with capacities induced by that subset. Unlike MoE-I\textsuperscript{2}'s staged layer/block selection, it adjusts full-model capacities and identities together. Unlike fixed-ranking allocation search in EvoESAP, it can change which experts occupy a given layer capacity. Equal-capacity comparisons and search-process analysis test the value of these freedoms.

\section{Method}
\label{sec:3}
\subsection{Joint representation under an exact global budget}
\label{sec:3.1}
For a model with $L$ MoE layers and $E_l$ routed experts in layer $l$, a binary mask $M$ selects experts: $M_{l,e}=1$ retains expert $e$. The layer capacity is
\[
k_l=\sum_e M_{l,e}.
\]
Every candidate satisfies
\[
\sum_{l=1}^{L}k_l=B,\qquad
k_l^{\min}\leq k_l\leq k_l^{\max}.
\]
The global budget $B$ is an equality constraint. Increasing one layer's capacity requires an equal reduction elsewhere. In each model studied here, routed experts have identical structures, so equal expert counts imply equal theoretical expert-parameter budgets.

The same mask determines both the number and identities of retained experts. Expert weights, router weights, and the original top-$k$ setting remain fixed. Routing is restricted to retained experts; shared experts and other dense modules are preserved.

\subsection{Importance-guided initialization}
\label{sec:3.2}
Importance scores construct initial expert sets. For example, REAP scores the conditional average of gate-weighted expert-output norms:
\[
S_{l,e}=\mathbb{E}\left[
g_{l,e}(x)\lVert f_{l,e}(x)\rVert_2
\mid e\text{ is selected}\right].
\]
TAPE initializes from sets produced by different scoring methods and previously available feasible sets, then adds local perturbations at different scales. Subsequent candidates can change both expert identities and layer capacities under the same global budget.

\subsection{Budget-preserving evolutionary operators}
\label{sec:3.3}
Represent a retained set as $S=\{(l,e):M_{l,e}=1\}$, where each element contains a layer and expert index and $|S|=B$.

An identity swap replaces one expert within the same layer. For $(l,e^-)\in S$ and $(l,e^+)\notin S$,
\[
S'=\bigl(S\setminus\{(l,e^-)\}\bigr)\cup\{(l,e^+)\}.
\]
All layer capacities remain unchanged: $k'_j=k_j$.

A capacity transfer removes an expert from one layer and adds one to another. For $l\neq m$, $(l,e^-)\in S$, and $(m,f^+)\notin S$,
\[
S'=\bigl(S\setminus\{(l,e^-)\}\bigr)\cup\{(m,f^+)\},
\qquad
k'_l=k_l-1,\quad k'_m=k_m+1.
\]
The operation requires $k_l>k_l^{\min}$ and $k_m<k_m^{\max}$. Other layer capacities are unchanged, and the total remains $B$.

Set crossover recombines parents $S_a,S_b$. Let $C=S_a\cap S_b$ and $D=S_a\triangle S_b$. Because both parents contain $B$ experts, $|D|=2(B-|C|)$. Randomly choose $|D|/2$ elements of $D$ as $T$, producing complementary children:
\[
\widetilde S_1=C\cup T,\qquad
\widetilde S_2=C\cup(D\setminus T).
\]
Both satisfy $|\widetilde S_1|=|\widetilde S_2|=B$. If a layer exceeds its allowed range, cross-layer transfers repair its capacity without changing the total: $S'_j=\mathcal R(\widetilde S_j)$.

Joint mutation combines transfers with identity swaps. NSGA-II uses non-dominated sorting and crowding distance to retain candidates with different trade-offs between the two tasks.

\subsection{Task-aware objectives}
\label{sec:3.4}
Search uses labeled questions from the target tasks' training splits without updating weights. For question $x$ and answer option $a_j$, compute
\[
s_M(x,a_j)=\frac{1}{c(a_j)}\sum_{u=1}^{|a_j|}
\log p_M(a_{j,u}\mid x,a_{j,<u}).
\]
Setting $c(a_j)=1$ gives total option log likelihood; setting it to the option's character count gives length-normalized scores. Dividing by temperature $\tau$ and applying softmax produces the candidate answer distribution $q_M$. The original model produces $p_0$ in the same way. For task $t$,
\[
L_t(M)=\mathbb{E}_{(x,y)\in\mathcal D_t}
\left[\lambda_{\mathrm{CE}}\bigl(-\log q_M(y\mid x)\bigr)
+\lambda_{\mathrm{KL}}\,\mathrm{KL}\bigl(p_0(\cdot\mid x)\Vert q_M(\cdot\mid x)\bigr)\right].
\]
The non-negative weights sum to one. Correct-answer cross-entropy measures task prediction, while teacher KL constrains deviation from the original model's distribution over the answer options. All scores come from forward inference.

Evolution minimizes two separate task objectives:
\[
\min_M\bigl(L_{\mathrm{ARC}}(M),L_{\mathrm{OBQA}}(M)\bigr).
\]
Keeping them separate preserves task-specific trade-offs. Continuous losses provide finer discrimination than small-sample accuracy when candidates answer the same number of questions correctly.

\section{Experimental setup}
\label{sec:4}
\subsection{Models and compression budgets}
\label{sec:4.1}
Our main model is ERNIE-4.5-21B-A3B-PT, with Granite-3.1-1B-A400M-Base as a smaller-model evaluation. Their original top-6 and top-8 routing settings are preserved. Expert weights, router-scoring weights, and non-pruned parameters remain fixed.

ERNIE has 27 MoE layers with 64 routed experts each; Granite has 24 with 32 each.\footnote{See the published \href{https://huggingface.co/baidu/ERNIE-4.5-21B-A3B-PT/blob/main/config.json}{ERNIE configuration} and \href{https://huggingface.co/ibm-granite/granite-3.1-1b-a400m-base/blob/main/config.json}{Granite configuration}.} Expert retention refers only to routed experts. Theoretical total parameters assume removal of unretained experts and their router-output rows, with dense modules and shared experts preserved. Total-model parameter reduction is therefore smaller than the fraction of removed routed experts.

\begin{table}[htbp]
\centering\small
\caption{Model sizes and expert-retention budgets.}
\label{tab:budget}
\begin{tabularx}{\linewidth}{@{}lYYYYY@{}}
\toprule
Model & \shortstack{Expert\\retention} & \shortstack{Experts\\retained} & \shortstack{Total\\parameters (B)} & \shortstack{Parameter\\reduction} & \shortstack{Experts\\per layer} \\
\midrule
ERNIE & 100\% & 1,728 & 21.949 & 0.00\% & 64 \\
ERNIE & 75\% & 1,296 & 16.851 & 23.22\% & 44--52 \\
ERNIE & 50\% & 864 & 11.754 & 46.45\% & 28--36 \\
Granite & 100\% & 768 & 1.335 & 0.00\% & 32 \\
Granite & 75\% & 576 & 1.033 & 22.63\% & 20--28 \\
Granite & 50\% & 384 & 0.731 & 45.25\% & 12--20 \\
\bottomrule
\end{tabularx}
\end{table}

All methods retain the same total number of experts for a given model and retention level. Scoring baselines use uniform layer capacities; TAPE redistributes those slots within the specified bounds.

\subsection{Tasks, data, and metrics}
\label{sec:4.2}
We evaluate ARC-Easy and OpenBookQA multiple-choice question answering. Small subsets of their official training splits guide search and final selection. The complete official test splits contain 2,376 ARC-Easy questions and 500 OpenBookQA questions.

The primary metric is character-length-normalized accuracy, \texttt{acc\_norm}. Each option's conditional log likelihood is divided by its character count, and the highest-scoring option is selected. This is the metric field selected in \href{https://arxiv.org/html/2510.13999v3#A3}{REAP Appendix C}~\citep{ref1}.
\[
\mathrm{Avg}=\frac{\mathrm{Acc}_{\mathrm{ARC}}+\mathrm{Acc}_{\mathrm{OBQA}}}{2}.
\]
Performance retention compares this equal-weight mean to the original model:
\[
\mathrm{Ret}=
\frac{\mathrm{Avg}_{\mathrm{pruned}}}{\mathrm{Avg}_{\mathrm{original}}}\times100\%.
\]
For example, an original average of 60\% and a pruned average of 54\% give 90\% performance retention. Tables report percentages. Scores and differences are calculated before rounding to two decimal places.

\subsection{Baselines}
\label{sec:4.3}
\textbf{Original} is the unpruned model. \textbf{REAP} ranks experts by gate-weighted output norm and retains the highest-scoring experts within each layer. \textbf{Frequency} ranks them by routing counts on calibration data.

\textbf{EAN-Mean} uses mean output norm conditional on activation, without gate weighting. It is distinguished from the accumulated-norm EAN variant in \href{https://arxiv.org/html/2510.13999v3#A4}{REAP Appendix D}~\citep{ref1}. ERNIE includes this baseline separately. \textbf{REAP+Frequency} sums standardized within-layer REAP scores and standardized log-frequency scores. Every baseline has the same total expert budget as TAPE.

\subsection{Search settings and reproducibility}
\label{sec:4.4}
Both models and retention levels use the same search and selection procedure. Each task provides 64 training questions for search and a disjoint 256 for final selection. NSGA-II uses a population of 16, 16 offspring per generation, and 15 generations, evaluating 256 distinct candidates per experiment.

All experiments use one NVIDIA A100-SXM4-80GB with BF16 forward inference. Granite evaluates two model replicas on the same GPU; ERNIE uses one. Recorded search-stage runtimes are 3.88--4.93 minutes for Granite and 75.01 minutes for ERNIE at 50\% retention. Detailed timing definitions and resource measurements are reported separately in \texttt{docs/COMPUTE\_COSTS.md}. Code, fixed inputs, and reproduction commands are available at \url{https://github.com/YuqiHuang2003/Tape}.

Both tasks use character-normalized option scores, temperature $\tau=0.25$, and weights $\lambda_{\mathrm{CE}}=0.75$, $\lambda_{\mathrm{KL}}=0.25$. Initial populations include REAP, Frequency, REAP+Frequency, a historical feasible TAPE set for the corresponding model and budget, and local neighbors.

After search, final selection considers all 16 initial candidates and eight generated candidates selected by search-set mean \texttt{acc\_norm}. The additional training subset selects the output by mean \texttt{acc\_norm}; ties are broken by the lower single-task accuracy, normalized label cross-entropy, and a fixed candidate identifier. The expert set is frozen before test evaluation. ERNIE 75\% identity-only and allocation-only ablations use the same sample counts, candidate budget, and selection rule, with initialization constrained to the respective search variables.

\section{Results and analysis}
\label{sec:5}
\subsection{ERNIE at 75\% and 50\% expert retention}
\label{sec:5.1}
Table~\ref{tab:ernie} reports full-test results. TAPE outputs are selected using the training-based rule in Section~\ref{sec:4.4}.

\begin{table}[htbp]
\centering\small
\caption{ERNIE results at two expert-retention budgets.}
\label{tab:ernie}
\begin{tabularx}{\linewidth}{@{}lXrrrr@{}}
\toprule
\shortstack{Expert\\retention} & Method & \shortstack{ARC-Easy\\Test (\%)} & \shortstack{OpenBookQA\\Test (\%)} & Avg (\%) & \shortstack{Performance\\retained} \\
\midrule
100\% & Original & 78.24 & 46.20 & 62.22 & 100.00\% \\
\midrule
75\% & REAP & 75.80 & 41.20 & 58.50 & 94.02\% \\
75\% & Frequency & 71.38 & 37.20 & 54.29 & 87.25\% \\
75\% & EAN-Mean & 74.07 & 42.20 & 58.14 & 93.44\% \\
75\% & REAP+Frequency & 73.23 & 38.40 & 55.82 & 89.71\% \\
75\% & TAPE & 80.93 & 44.40 & 62.67 & 100.72\% \\
\midrule
50\% & REAP & 63.34 & 34.20 & 48.77 & 78.38\% \\
50\% & Frequency & 59.30 & 29.20 & 44.25 & 71.12\% \\
50\% & EAN-Mean & 67.42 & 32.60 & 50.01 & 80.38\% \\
50\% & REAP+Frequency & 61.95 & 34.00 & 47.98 & 77.11\% \\
50\% & TAPE & 68.52 & 33.00 & 50.76 & 81.58\% \\
\bottomrule
\end{tabularx}
\end{table}

At 75\% retention, TAPE reaches 62.67\% Avg, exceeding REAP by 4.17 percentage points and EAN-Mean by 4.53. Compared with REAP, ARC-Easy and OpenBookQA improve by 5.13 and 3.20 points. Theoretical total parameters decrease by 23.22\%, while Avg exceeds the original model by 0.45 points, retaining 100.72\% of its average performance. Relative to the original model, ARC-Easy improves by 2.69 points and OpenBookQA decreases by 1.80 points.

At 50\% retention, TAPE achieves 50.76\% Avg, 1.99 points above REAP and 0.75 above EAN-Mean. It answers 26 more ARC-Easy questions and two more OpenBookQA questions correctly than EAN-Mean. With a 46.45\% reduction in theoretical total parameters, performance retention is 81.58\%. Section~\ref{sec:5.4} examines continued evolution at 75\% retention.

\subsection{Granite: compression of the smaller model}
\label{sec:5.2}
\begin{table}[htbp]
\centering\small
\caption{Granite results at two expert-retention budgets.}
\label{tab:granite}
\begin{tabularx}{\linewidth}{@{}lXrrrr@{}}
\toprule
\shortstack{Expert\\retention} & Method & \shortstack{ARC-Easy\\Test (\%)} & \shortstack{OpenBookQA\\Test (\%)} & Avg (\%) & \shortstack{Performance\\retained} \\
\midrule
100\% & Original & 67.59 & 38.00 & 52.80 & 100.00\% \\
\midrule
75\% & REAP & 46.89 & 30.60 & 38.74 & 73.38\% \\
75\% & Frequency & 52.69 & 30.40 & 41.55 & 78.69\% \\
75\% & REAP+Frequency & 53.45 & 32.80 & 43.13 & 81.68\% \\
75\% & TAPE & 54.12 & 34.40 & 44.26 & 83.84\% \\
\midrule
50\% & REAP & 33.46 & 26.80 & 30.13 & 57.07\% \\
50\% & Frequency & 28.49 & 29.00 & 28.75 & 54.45\% \\
50\% & REAP+Frequency & 35.90 & 28.00 & 31.95 & 60.52\% \\
50\% & TAPE & 37.67 & 27.00 & 32.33 & 61.24\% \\
\bottomrule
\end{tabularx}
\end{table}

At 75\% retention, TAPE reaches 44.26\% Avg, outperforming all direct-scoring methods: 5.52 points above REAP and 1.14 above REAP+Frequency. Performance retention is 83.84\%.

At 50\%, TAPE reaches 32.33\% Avg, exceeding REAP by 2.20 points and REAP+Frequency by 0.38, with 61.24\% performance retention. Compared with combined scoring, ARC-Easy improves by 1.77 points while OpenBookQA decreases by 1.00. The average improvement therefore comes from ARC-Easy rather than a gain shared by both tasks.

The two models share sample counts, search budgets, objectives, and selection rules. ERNIE retains a larger fraction of its original performance at both budgets. On Granite, performance retention decreases from 83.84\% to 61.24\% when expert retention falls from 75\% to 50\%, showing a substantial quality cost under stronger compression.

\subsection{Ablations}
\label{sec:5.3}
Ablations use ERNIE at 75\% expert retention. \emph{w/o Layer Allocation} fixes every layer at 48 experts and searches identities only. \emph{w/o Identity Search} searches capacities while selecting experts by fixed REAP rankings. These are independent searches with the same sample counts, candidate budget, and final selection rule. \emph{w/o Evolved Identities} fixes the full TAPE output's capacities and replaces its expert set with REAP-ranked choices.

\begin{table}[htbp]
\centering\small
\caption{Ablations on ERNIE at 75\% expert retention.}
\label{tab:ablation}
\begin{tabularx}{\linewidth}{@{}Xrrr@{}}
\toprule
Method & ARC-Easy (\%) & OpenBookQA (\%) & Avg (\%) \\
\midrule
TAPE & 80.93 & 44.40 & 62.67 \\
w/o Layer Allocation & 77.99 & 41.60 & 59.79 \\
w/o Identity Search & 76.43 & 40.40 & 58.42 \\
w/o Evolved Identities & 76.14 & 40.00 & 58.07 \\
\bottomrule
\end{tabularx}
\end{table}

TAPE exceeds w/o Layer Allocation by 2.87 points and w/o Identity Search by 4.25. The joint search produces a higher-scoring set than either single-variable search.

With TAPE's layer capacities fixed, REAP-ranked selection reaches 58.07\% Avg, 4.60 points below TAPE. The expert count and every layer capacity are identical; only expert identities differ. Changing uniform REAP capacities to TAPE's capacities moves Avg from 58.50\% to 58.07\%, while using TAPE's identities at those capacities raises it to 62.67\%. Task-level evaluation improves the expert combination beyond the fixed importance ordering.

\subsection{ERNIE 75\%: evolution and expert-set changes}
\label{sec:5.4}
We compare TAPE after initialization (generation 0) with its full 15-generation run. The generation-0 control selects from all 16 initial candidates; the full run also considers the generated candidates specified in Section~\ref{sec:4.4}. Both use the same additional training questions and final selection rule, followed by full-test evaluation. The full-run row reuses the 75\% main result from Section~\ref{sec:5.1}.

\begin{table}[htbp]
\centering\small
\caption{Initialization-only and full-search results on ERNIE at 75\% expert retention.}
\label{tab:evolution}
\begin{tabularx}{\linewidth}{@{}Xrrr@{}}
\toprule
Method & ARC-Easy (\%) & OpenBookQA (\%) & Avg (\%) \\
\midrule
TAPE (0 generations) & 80.60 & 41.60 & 61.10 \\
TAPE (15 generations) & 80.93 & 44.40 & 62.67 \\
\bottomrule
\end{tabularx}
\end{table}

The selected full-search candidate was generated in generation 14. Relative to the generation-0 output, it changes the capacities of 16 layers and replaces 26 experts. Avg increases from 61.10\% to 62.67\%, a 1.57-point gain. ARC-Easy improves by 0.34 points and OpenBookQA by 2.80, making the latter the main source of improvement. Continued evolution finds a combination outside the initial population that is selected on training questions and performs better on test.

\section{Expert-selection behavior}
\label{sec:6}
We examine how TAPE's retained experts differ from importance-ranked sets and how those choices associate with task losses.

\subsection{Retention decisions relative to REAP}
\label{sec:6.1}
ERNIE has 1,728 layer-indexed routed experts. At 75\% retention, 1,296 are kept and 432 removed. We first compare TAPE with the uniform-capacity REAP reference in the initial population.

\begin{table}[htbp]
\centering\small
\caption{Expert-retention decisions of TAPE and uniform-capacity REAP on ERNIE 75\%.}
\label{tab:decisions}
\begin{tabularx}{\linewidth}{@{}Xr@{}}
\toprule
Decision relationship & Experts \\
\midrule
Retained by both & 1179 \\
Removed by both & 315 \\
Removed by REAP, retained by TAPE & 117 \\
Retained by REAP, removed by TAPE & 117 \\
\bottomrule
\end{tabularx}
\end{table}

Of the experts removed by TAPE, 72.9\% are also removed by REAP. Differences comprise 117 experts restored relative to REAP and 117 newly removed experts.

To separate capacity and identity changes, we construct another REAP set using TAPE's exact layer capacities. Relative to this reference, TAPE retains 128 experts REAP would remove and removes 128 REAP would retain, spanning all 27 MoE layers. Their test averages are 62.67\% and 58.07\%. The counts 117 and 128 refer to different reference sets: uniform-capacity and capacity-matched REAP, respectively.

\subsection{Associations between expert retention and task performance}
\label{sec:6.2}
For each expert, we partition the 256 ERNIE 75\% search candidates by whether they retain it, then compare task objectives. We standardize objectives and linearly adjust for generation and layer capacities. Results are reported only when both groups have at least 26 candidates. Retention association is the mean adjusted objective in the removal group minus that in the retention group. Positive values associate retention with lower loss; negative values associate it with higher loss.

Among the 128 restored experts in the capacity-matched comparison, 23 meet the sample requirement; 10 have a retention-favorable mean association across the tasks (43.5\%). Among the 128 newly removed experts, 35 qualify and 14 favor retention (40.0\%). These proportions are close and show no clear population-level separation.

\begin{table}[htbp]
\centering\small
\caption{Selected examples of task-level retention associations at matched capacities.}
\label{tab:association}
\begin{tabularx}{\linewidth}{@{}lrrlXX@{}}
\toprule
Expert & \shortstack{Layer\\capacity} & \shortstack{REAP\\rank} & Decision & \shortstack{ARC\\association} & \shortstack{OBQA\\association} \\
\midrule
L11-E20 & 46 & 50 / 64 & Retain & Favorable\newline(+0.113) & Favorable\newline(+0.147) \\
L21-E46 & 51 & 53 / 64 & Retain & Favorable\newline(+0.100) & Favorable\newline(+0.085) \\
L21-E59 & 51 & 50 / 64 & Remove & Unfavorable\newline(\ensuremath{-}0.176) & Unfavorable\newline(\ensuremath{-}0.156) \\
L25-E15 & 49 & 13 / 64 & Remove & Unfavorable\newline(\ensuremath{-}0.103) & Unfavorable\newline(\ensuremath{-}0.157) \\
\bottomrule
\end{tabularx}
\end{table}

The table lists two restored experts with high mean retention association and two newly removed experts with low mean association. For layer 21 at capacity 51, REAP retains rank-50 E59 and removes rank-53 E46; TAPE does the reverse. E46 has favorable retention associations on both tasks, while E59 has unfavorable ones. This example illustrates a difference between importance order and task-associated behavior. TAPE evaluates complete pruned combinations instead of limiting selection to a fixed ordering of individual scores.

\section{Conclusion}
\label{sec:7}
TAPE jointly searches layer capacities and expert identities under a fixed global expert budget, using task-level multi-objective evaluation without updating model weights. Its mean test accuracy exceeds the direct-scoring baselines at both budgets on both models. On Granite, it improves over combined scoring by 1.14 points at 75\% retention and 0.38 at 50\%.

ERNIE ablations and capacity-matched comparisons support joint optimization. At 75\% retention, continued evolution raises Avg from 61.10\% to 62.67\%. Expert-selection analysis identifies cases where importance rankings differ from retention-associated task losses, showing how task-level evaluation can complement individual importance scores.

\section{Limitations and future work}
\label{sec:8}
Evaluation covers ARC-Easy and OpenBookQA, with small labeled subsets of these tasks used for search and selection. Future work can extend to more question-answering and generation tasks, evaluate tasks excluded from search, and study transfer of task-aware expert selection.

Exploratory layer-capacity analysis did not identify a clear, stable depth preference. The task-specific distribution of useful experts remains unresolved. Holding capacities and other expert choices fixed while swapping individual experts, and relating the resulting score changes to routing distributions and input types, could clarify how combinations of experts contribute to different tasks.

\begin{thebibliography}{99}
\item[] Twenty studies from 2024--2026 are listed below; preprints are identified separately.

\bibitem{ref1}
Lasby et al. \href{https://proceedings.iclr.cc/paper_files/paper/2026/hash/ed93b2b5722acc2341d421b8916404a1-Abstract-Conference.html}{REAP the Experts: Why Pruning Prevails for One-Shot MoE compression}. ICLR, 2026.

\bibitem{ref2}
Lu et al. \href{https://aclanthology.org/2024.acl-long.334/}{Not All Experts are Equal: Efficient Expert Pruning and Skipping for Mixture-of-Experts Large Language Models}. ACL, 2024.

\bibitem{ref3}
Chen et al. \href{https://aclanthology.org/2025.acl-long.633/}{EAC-MoE: Expert-Selection Aware Compressor for Mixture-of-Experts Large Language Models}. ACL, 2025.

\bibitem{ref4}
Lee et al. \href{https://aclanthology.org/2025.acl-long.671/}{STUN: Structured-Then-Unstructured Pruning for Scalable MoE Pruning}. ACL, 2025.

\bibitem{ref5}
Li et al. \href{https://proceedings.iclr.cc/paper_files/paper/2024/hash/3d09a88c3372cdb79401fde592ca4db8-Abstract-Conference.html}{Merge, Then Compress: Demystify Efficient SMoE with Hints from Its Routing Policy}. ICLR, 2024.

\bibitem{ref6}
Chen et al. \href{https://proceedings.mlr.press/v267/chen25aq.html}{Retraining-free Merging of Sparse MoE via Hierarchical Clustering}. ICML, 2025.

\bibitem{ref7}
Zhang et al. \href{https://aclanthology.org/2025.findings-acl.4/}{Diversifying the Expert Knowledge for Task-Agnostic Pruning in Sparse Mixture-of-Experts}. Findings of ACL, 2025.

\bibitem{ref8}
Zhou et al. \href{https://aclanthology.org/2025.findings-emnlp.820/}{Dropping Experts, Recombining Neurons: Retraining-Free Pruning for Sparse Mixture-of-Experts LLMs}. Findings of EMNLP, 2025.

\bibitem{ref9}
Nguyen et al. \href{https://proceedings.iclr.cc/paper_files/paper/2025/hash/341d64b067e06c6ebea97a498c90d598-Abstract-Conference.html}{CAMEx: Curvature-aware Merging of Experts}. ICLR, 2025.

\bibitem{ref10}
Nguyen et al. \href{https://proceedings.iclr.cc/paper_files/paper/2026/hash/ef9f16cd0cb8d7797ea35c0ee7964cdf-Abstract-Conference.html}{Expert Merging in Sparse Mixture of Experts with Nash Bargaining}. ICLR, 2026.

\bibitem{ref11}
Li et al. \href{https://proceedings.mlr.press/v267/li25az.html}{MoE-SVD: Structured Mixture-of-Experts LLMs Compression via Singular Value Decomposition}. ICML, 2025.

\bibitem{ref12}
Gu et al. \href{https://proceedings.mlr.press/v267/gu25c.html}{Delta Decompression for MoE-based LLMs Compression}. ICML, 2025.

\bibitem{ref13}
Li et al. \href{https://ojs.aaai.org/index.php/AAAI/article/view/39464}{Sub-MoE: Efficient Mixture-of-Expert LLMs Compression via Subspace Expert Merging}. AAAI, 2026.

\bibitem{ref14}
Huang et al. \href{https://proceedings.iclr.cc/paper_files/paper/2025/hash/abc1943857a42935ceacff03c524bb44-Abstract-Conference.html}{Mixture Compressor for Mixture-of-Experts LLMs Gains More}. ICLR, 2025.

\bibitem{ref15}
Liu et al. \href{https://arxiv.org/abs/2407.00945}{Efficient Expert Pruning for Sparse Mixture-of-Experts Language Models: Enhancing Performance and Reducing Inference Costs}. arXiv preprint, 2024.

\bibitem{ref16}
Yang et al. \href{https://aclanthology.org/2024.findings-emnlp.612/}{MoE-I\textsuperscript{2}: Compressing Mixture of Experts Models through Inter-Expert Pruning and Intra-Expert Low-Rank Decomposition}. Findings of EMNLP, 2024.

\bibitem{ref17}
Liu et al. \href{https://arxiv.org/abs/2603.06003}{EvoESAP: Non-Uniform Expert Pruning for Sparse MoE}. arXiv preprint, 2026.

\bibitem{ref18}
Akiba et al. \href{https://www.nature.com/articles/s42256-024-00975-8}{Evolutionary Optimization of Model Merging Recipes}. Nature Machine Intelligence, 2025.

\bibitem{ref19}
Kuroki et al. \href{https://proceedings.iclr.cc/paper_files/paper/2025/hash/755acd0c7c07180d78959b6d89768207-Abstract-Conference.html}{Agent Skill Acquisition for Large Language Models via CycleQD}. ICLR, 2025.

\bibitem{ref20}
Mencattini et al. \href{https://proceedings.mlr.press/v267/mencattini25a.html}{MERGE\textsuperscript{3}: Efficient Evolutionary Merging on Consumer-grade GPUs}. ICML, 2025.

\end{thebibliography}
\end{document}
